\chapter{Splenectomy}

\section*{What Is This Surgery?}
A splenectomy is the surgical removal of the spleen, a vital lymphoid organ situated in the left upper quadrant of the abdomen, beneath the diaphragm and posterior to the stomach. This procedure is performed for a diverse range of clinical scenarios, including acute trauma, various hematologic disorders, and certain oncologic conditions. The spleen plays crucial roles in the body's immune system, filtering blood, removing old or damaged red blood cells, and storing platelets. Its removal can significantly impact a patient's hematologic profile and immune function, necessitating careful consideration and comprehensive postoperative management to mitigate the lifelong risks, particularly overwhelming post-splenectomy infection (OPSI). The surgery can be performed via an open or minimally invasive (laparoscopic) approach, with the choice depending on the clinical context and splenic characteristics.

\section*{Alternative Names}
\begin{itemize}
  \item Spleen removal surgery
  \item Splenic resection
  \item Total splenectomy
  \item Partial splenectomy (when only a portion of the spleen is removed)
  \item Laparoscopic splenectomy (keyhole spleen surgery)
  \item Open splenectomy (traditional spleen surgery)
\end{itemize}

\section*{Relevant Surgical Anatomy}
The spleen is an ovoid, reddish-purple organ, typically measuring 12 cm in length, 7 cm in width, and 3-4 cm in thickness, weighing approximately 150-200 grams in adults. It is located in the left upper quadrant of the abdomen, deep to the 9th, 10th, and 11th ribs, nestled between the fundus of the stomach and the left hemidiaphragm. Its long axis lies along the 10th rib. Key anatomical relations include the diaphragm superiorly, the greater curvature of the stomach anteriorly, the tail of the pancreas medially, the left kidney inferiorly, and the left colic flexure inferiorly and posteriorly.

The spleen is suspended by several peritoneal ligaments. The gastrosplenic ligament connects the greater curvature of the stomach to the splenic hilum and contains the short gastric arteries and veins, which supply the fundus of the stomach. The splenorenal (or lienorenal) ligament extends from the splenic hilum to the posterior abdominal wall over the left kidney and contains the main splenic artery and vein, as well as the tail of the pancreas. Other ligaments include the splenocolic ligament (connecting the spleen to the left colic flexure) and the phrenicocolic ligament (supporting the left colic flexure, indirectly supporting the spleen). The splenic artery, a branch of the celiac trunk, is tortuous and runs along the superior border of the pancreas, giving off pancreatic branches and short gastric arteries before entering the splenic hilum. The splenic vein drains the spleen and runs posterior to the pancreas, joining the superior mesenteric vein to form the portal vein. Lymphatic drainage follows the splenic vessels to pancreaticosplenic lymph nodes and then to celiac nodes. The spleen is innervated by sympathetic fibers from the celiac plexus. Accessory spleens (splenunculi) are present in approximately 10-30\% of individuals, most commonly found near the splenic hilum, along the gastrosplenic ligament, or in the omentum, and are surgically significant as they can lead to recurrence of hematologic conditions if not removed.

\section{Surgical Principle \& Rationale}
The fundamental principle of splenectomy is to eliminate a diseased or damaged spleen that is either causing life-threatening hemorrhage, contributing to the destruction or sequestration of blood components, or acting as a primary site of pathology. The spleen's functions include filtering senescent red blood cells, sequestering platelets, participating in immune surveillance by producing antibodies, and removing encapsulated bacteria. When these functions become dysregulated, as in hypersplenism or autoimmune cytopenias, or when the organ is acutely injured, its removal becomes a therapeutic necessity.

The surgical rationale is to achieve complete removal of the splenic parenchyma while meticulously controlling the splenic artery and vein to prevent hemorrhage. Technical goals include precise dissection to avoid injury to adjacent organs such as the tail of the pancreas, the stomach, and the left hemidiaphragm. For certain hematologic conditions, the goal is to remove the primary site of blood cell destruction or sequestration, thereby allowing peripheral blood counts to normalize. In cases of trauma, the primary goal is to achieve immediate hemostasis and prevent exsanguination. For oncologic indications, the aim is to resect the primary tumor or reduce tumor burden. The identification and removal of any accessory spleens are also critical, particularly in cases of hematologic disorders, to prevent disease recurrence.

\section{History \& Surgical Evolution}
The concept of splenectomy dates back centuries, with early attempts often resulting in fatal outcomes due to uncontrolled hemorrhage. The first documented successful splenectomy is often attributed to Adrian Zaccarelli in 1549, though details are scarce and survival was rare. More reliably, the first successful elective splenectomy for a hematologic disorder (leukemia) was performed by Franz Christian von Hildenbrand in 1826. However, it was not until the late 19th and early 20th centuries that surgical techniques and anesthetic advancements made splenectomy a more viable procedure.

Key milestones in the evolution of splenectomy include:
\begin{itemize}
  \item \textbf{1887:} Theodor Billroth performed a successful splenectomy for a splenic tumor, marking a significant step in the procedure's acceptance for non-traumatic indications.
  \item \textbf{Early 20th Century:} The procedure gained wider acceptance for trauma and various hematological conditions, such as hereditary spherocytosis and immune thrombocytopenia (ITP), as understanding of splenic physiology improved.
  \item \textbf{1950s-1970s:} Increased recognition of the risk of overwhelming post-splenectomy infection (OPSI) led to more cautious indications and the development of prophylactic vaccination strategies.
  \item \textbf{1991:} Delaitre and Maignien performed the first successful laparoscopic splenectomy, which revolutionized the approach by offering reduced invasiveness, shorter hospital stays, and quicker recovery times for elective cases. This minimally invasive technique has since become the preferred approach for most elective splenectomies.
\end{itemize}
The evolution continues with advancements in non-operative management for splenic trauma and refined medical therapies for hematologic disorders, further refining the indications for surgical intervention.

\section*{Clinical Indications}
Splenectomy is indicated for a variety of conditions, broadly categorized into trauma, hematologic disorders, and oncologic or other pathologies.

\begin{itemize}
  \item Traumatic splenic rupture with hemodynamic instability or significant ongoing hemorrhage that is refractory to non-operative management.
  \item Immune thrombocytopenia (ITP) refractory to first-line medical therapy (e.g., corticosteroids, IVIG) and causing severe bleeding or a high risk of hemorrhage.
  \item Hereditary spherocytosis or other congenital hemolytic anemias (e.g., elliptocytosis, pyruvate kinase deficiency) leading to severe anemia, transfusion dependence, or symptomatic splenomegaly.
  \item Hypersplenism causing significant pancytopenia, leukopenia, or thrombocytopenia, often associated with portal hypertension or myeloproliferative disorders.
  \item Splenic abscess unresponsive to antibiotic therapy and percutaneous drainage, or in cases of multiloculated abscesses.
  \item Primary splenic tumors, such as lymphoma, hemangioma, or angiosarcoma, or metastatic disease to the spleen, particularly when localized or causing symptoms.
  \item Splenic vein thrombosis with associated portal hypertension and recurrent variceal bleeding, especially when isolated to the splenic vein.
  \item Massive splenomegaly causing severe pain, early satiety due to mass effect, or significant cytopenias, often seen in myelofibrosis or chronic leukemias.
  \item As part of a staging procedure for certain lymphomas (though less common now with advanced imaging) or as part of a larger resection for adjacent organ malignancies (e.g., distal pancreatectomy with splenectomy).
\end{itemize}

\section*{Clinical Positioning}
Splenectomy can serve as both a primary and a secondary intervention, depending on the clinical context and the underlying pathology.

In cases of acute, life-threatening traumatic splenic rupture with hemodynamic instability, splenectomy is often a primary, emergent procedure aimed at achieving immediate hemostasis and saving the patient's life. While non-operative management (NOM) is increasingly favored for hemodynamically stable splenic trauma, surgical intervention becomes primary when NOM fails or is contraindicated.

Conversely, for many chronic hematologic conditions like immune thrombocytopenia (ITP) or hereditary spherocytosis, splenectomy is typically a secondary or salvage procedure, considered only after medical management has failed or proven inadequate. For instance, in ITP, corticosteroids, intravenous immunoglobulin (IVIG), and thrombopoietin receptor agonists are usually first-line therapies, with splenectomy reserved for patients who are refractory, relapse, or experience severe side effects from medical treatments. Similarly, for hereditary spherocytosis, splenectomy is usually deferred until after childhood to minimize the risk of OPSI, and only performed if the anemia is severe or transfusion-dependent. For splenic tumors, it may be a primary treatment if the tumor is localized and resectable, or a secondary procedure if it is part of a broader oncologic management strategy, such as debulking or as an adjunct to chemotherapy. The increasing success of non-operative management for splenic trauma and advancements in medical therapies for hematologic disorders have generally shifted splenectomy towards a more secondary role in elective settings.

\section*{Surgical Technique \& Procedure}
Splenectomy is performed under general anesthesia and can be approached either openly or laparoscopically. The choice of approach depends on the patient's condition, splenic size, presence of adhesions, and surgeon's expertise. Laparoscopic splenectomy is generally preferred for elective cases due to reduced pain, shorter hospital stay, and quicker recovery, provided the spleen is not excessively large or there are no extensive adhesions.

\textbf{Laparoscopic Splenectomy (Preferred for Elective Cases):}
\begin{itemize}
  \item \textbf{Anesthesia \& Positioning:} General anesthesia is induced. The patient is typically placed in a right lateral decubitus position, allowing gravity to retract the bowel away from the left upper quadrant. The operating table may be tilted.
  \item \textbf{Access \& Trocar Placement:} Pneumoperitoneum is established using a Veress needle or an open Hasson technique, inflating the abdomen with CO\textsubscript{2} to 12-15 mmHg. Typically, 3-4 trocars are placed in the left upper quadrant and epigastrium, allowing for optimal triangulation and instrument manipulation.
  \item \textbf{Exploration \& Mobilization:} The abdomen is explored to confirm the diagnosis and assess splenic size and surrounding anatomy. The spleen is mobilized by dividing its peritoneal attachments. The splenocolic and splenorenal ligaments are divided using an energy device (e.g., harmonic scalpel, LigaSure) to free the inferior and posterior aspects of the spleen.
  \item \textbf{Short Gastric Vessel Division:} The gastrosplenic ligament is carefully divided, ligating or sealing the short gastric vessels that connect the spleen to the greater curvature of the stomach. This step requires precision to avoid gastric injury.
  \item \textbf{Hilar Dissection \& Vessel Ligation:} The splenic artery and vein are identified at the hilum, often after retracting the tail of the pancreas inferiorly. The vessels are individually ligated and divided using clips, staplers, or energy devices. Extreme caution is exercised to protect the pancreatic tail from thermal or mechanical injury.
  \item \textbf{Accessory Spleen Search:} A thorough search for accessory spleens is performed, especially in cases of hematologic indications, in common locations such as the splenic hilum, gastrosplenic ligament, omentum, or pelvis. Any identified accessory spleens are removed.
  \item \textbf{Specimen Retrieval:} Once fully mobilized and devascularized, the spleen is placed into an endoscopic retrieval bag. The bag is then extracted through one of the port sites, often enlarged, or through a small Pfannenstiel incision. For larger spleens, morcellation within the bag may be necessary to facilitate removal.
  \item \textbf{Closure:} Hemostasis is confirmed, and the trocars are removed under direct vision. The port sites are closed in layers, with fascial closure for sites larger than 10 mm to prevent hernia. Skin incisions are closed with sutures or staples.

\end{itemize}

\textbf{Open Splenectomy (For Trauma, Massive Splenomegaly, or Complex Cases):}
\begin{itemize}
  \item \textbf{Anesthesia \& Positioning:} General anesthesia is induced, and the patient is typically positioned supine.
  \item \textbf{Incision:} A left subcostal (Kocher) incision or an upper midline incision is made, providing wide access to the left upper quadrant.
  \item \textbf{Mobilization:} The surgeon gains access to the abdominal cavity, identifies the spleen, and carefully mobilizes it by dividing its ligamentous attachments (gastrosplenic, splenorenal, splenocolic, phrenicocolic) using electrocautery or ligatures.
  \item \textbf{Hilar Dissection \& Vessel Ligation:} The splenic artery and vein are meticulously identified, ligated, and divided, usually at the hilum. Care is taken to avoid injury to the tail of the pancreas, which lies in close proximity.
  \item \textbf{Accessory Spleen Search:} A thorough search for accessory spleens is performed.
  \item \textbf{Removal \& Closure:} Once all attachments and vascular supplies are divided, the spleen is removed. Hemostasis is ensured, and the incision is closed in layers. A drain may be placed if there is concern for pancreatic injury or significant oozing.
\end{itemize}

\section*{Preoperative Workup \& Preparation}
Thorough preoperative preparation is crucial to optimize patient outcomes and minimize risks associated with splenectomy.

\begin{itemize}
  \item \textbf{Medical Evaluation:} A comprehensive history and physical examination are performed, along with a review of all medications, allergies, and comorbidities. Cardiac and anesthesia clearance are obtained, especially for patients with significant cardiovascular or pulmonary disease.
  \item \textbf{Laboratory Tests:} Routine labs include a complete blood count (CBC) with differential, coagulation profile (PT/INR, PTT), liver and renal function tests, electrolytes, and blood type and cross-match for potential transfusion. For hematologic indications, specific tests like peripheral blood smear, Coombs test, or flow cytometry may be required.
  \item \textbf{Imaging Studies:} An abdominal CT scan with intravenous contrast is often performed to assess splenic size, vascular anatomy, identify any accessory spleens, and evaluate for other intra-abdominal pathology. Ultrasound may also be used.
  \item \textbf{Medication Management:} Any anticoagulant or antiplatelet medications may need to be held or bridged according to established protocols to minimize bleeding risk. Immunosuppressants may require adjustment.
  \item \textbf{NPO Status:} Patients are instructed to be NPO (nil per os) for a specified period (typically 6-8 hours for solids, 2 hours for clear liquids) before surgery to prevent aspiration.
  \item \textbf{Vaccinations:} Preoperative vaccinations against encapsulated bacteria are mandatory for elective splenectomy and should ideally be administered at least 2 weeks prior to surgery to allow for adequate antibody development. These include:
    \begin{itemize}
      \item Pneumococcal vaccine (PCV13 and PPSV23)
      \item Meningococcal vaccine (quadrivalent ACWY and serogroup B)
      \item Haemophilus influenzae type B (Hib) vaccine
    \end{itemize}
  \item \textbf{Prophylactic Antibiotics:} Broad-spectrum intravenous antibiotics (e.g., cefazolin) are administered within 60 minutes prior to skin incision to reduce the risk of surgical site infection.
  \item \textbf{Informed Consent:} Detailed informed consent is obtained, explaining the procedure, potential benefits, risks (including OPSI), alternatives, and the possibility of conversion from laparoscopic to open approach.
\end{itemize}

\section*{Contraindications \& Precautions}
\textbf{Absolute Contraindications:}
\begin{itemize}
  \item Uncorrectable severe coagulopathy or active, uncontrolled bleeding diathesis, unless the splenectomy itself is indicated to treat the coagulopathy (e.g., severe ITP with life-threatening hemorrhage).
  \item Severe cardiopulmonary disease that precludes safe general anesthesia and the physiological stress of surgery.
  \item Patient refusal or inability to provide informed consent.
  \item Diffuse metastatic cancer where splenectomy offers no palliative benefit and carries significant risk.
\end{itemize}

\textbf{Relative Contraindications and Precautions:}
\begin{itemize}
  \item Severe portal hypertension with extensive portosystemic collateral circulation and large varices, which significantly increases the risk of intraoperative hemorrhage and postoperative portal vein thrombosis.
  \item Extensive intra-abdominal adhesions from previous surgeries, making dissection challenging and increasing the risk of organ injury, potentially favoring an open approach.
  \item Active systemic infection or sepsis, unless the spleen is the primary source of the infection (e.g., splenic abscess), as surgery can exacerbate systemic inflammation.
  \item Pregnancy, where the risks to both mother and fetus must be carefully weighed against the benefits of surgery; usually deferred until after delivery if possible.
  \item Extreme splenomegaly (e.g., spleen extending below the umbilicus or weighing >2 kg), which may necessitate an open approach, increase operative time, and heighten the risk of bleeding and pancreatic injury.
  \item Patient refusal of blood products, which can complicate the management of potential intraoperative or postoperative hemorrhage.
  \item Coexisting severe liver disease, which can impair coagulation and increase the risk of complications.
  \item Inadequate preoperative vaccination against encapsulated bacteria, especially for elective cases, increasing the lifelong risk of OPSI.
\end{itemize}

\section*{Complications \& Risks}
Splenectomy, like any major surgical procedure, carries inherent risks. These can be broadly categorized into general surgical complications and procedure-specific complications.

\textbf{General Surgical Risks:}
\begin{itemize}
  \item \textbf{Bleeding:} Intraoperative or postoperative hemorrhage, potentially requiring blood transfusions, reoperation, or interventional radiology embolization.
  \item \textbf{Infection:} Surgical site infection (wound infection, intra-abdominal abscess), pneumonia, or sepsis.
  \item \textbf{Anesthesia Complications:} Allergic reactions to medications, respiratory depression, cardiac arrhythmias, myocardial infarction, stroke, or aspiration pneumonia.
  \item \textbf{Deep Vein Thrombosis (DVT) and Pulmonary Embolism (PE):} Formation of blood clots in the legs or lungs, a risk increased by surgery and immobility.
\item \textbf{Injury to adjacent organs:} Unintended damage to bowel, liver, or major vessels during abdominal entry or dissection.
\end{itemize}

\textbf{Procedure-Specific Risks:}
\begin{itemize}
  \item \textbf{Overwhelming Post-Splenectomy Infection (OPSI):} A rare but life-threatening septicemia, particularly from encapsulated bacteria (e.g., Streptococcus pneumoniae, Neisseria meningitidis, Haemophilus influenzae type B), with a high mortality rate (50-70\%). This risk persists lifelong.
  \item \textbf{Pancreatic Injury:} Damage to the tail of the pancreas, which lies in close proximity to the splenic hilum, potentially leading to acute pancreatitis, pancreatic fistula, or pseudocyst formation.
  \item \textbf{Gastric Injury:} Accidental perforation or injury to the stomach, especially during division of the short gastric vessels within the gastrosplenic ligament.
  \item \textbf{Thrombocytosis and Thromboembolic Events:} A significant increase in platelet count (thrombocytosis) post-splenectomy can lead to an increased risk of deep vein thrombosis (DVT), pulmonary embolism (PE), or portal vein thrombosis (PVT), particularly in patients with underlying myeloproliferative disorders.
  \item \textbf{Left Lower Lobe Atelectasis or Pleural Effusion:} Due to irritation of the diaphragm, reduced lung expansion post-surgery, or fluid collection.
  \item \textbf{Accessory Spleen Missed:} If an accessory spleen is present and not removed, it can lead to recurrence of the underlying hematologic disease (e.g., persistent thrombocytopenia in ITP).
  \item \textbf{Injury to the Diaphragm:} Can occur during dissection, potentially leading to a diaphragmatic hernia or chronic pain.
  \item \textbf{Hernia at Incision Site:} More common with open splenectomy or at enlarged port sites in laparoscopic procedures.
  \item \textbf{Adhesions:} Formation of scar tissue within the abdomen, potentially causing future bowel obstruction.
\end{itemize}

\section*{Postoperative Recovery Timeline}
Immediately after splenectomy, the patient is transferred to the Post-Anesthesia Care Unit (PACU) for close monitoring of vital signs, pain levels, and signs of bleeding or other immediate complications. Early mobilization is encouraged to prevent pulmonary complications and deep vein thrombosis. Pain management is initiated, often with intravenous analgesics, transitioning to oral medications as tolerated.

During the first 24-48 hours, patients are monitored for signs of infection (fever, chills), bleeding (tachycardia, hypotension, decreasing hemoglobin), and pancreatic injury (elevated amylase/lipase, abdominal pain). A complete blood count is typically monitored, especially the platelet count, which often spikes post-splenectomy, necessitating consideration of antiplatelet therapy in high-risk patients. Diet is advanced gradually from clear liquids to a regular diet as bowel function returns. Most patients are discharged within 2-5 days following an uncomplicated laparoscopic splenectomy, or 5-7 days for an open procedure. Patients are advised on wound care, activity restrictions (avoiding heavy lifting or strenuous exercise for 4-6 weeks), and the importance of lifelong vigilance for signs of infection. Education on the symptoms of OPSI and the need for prompt medical attention for any fever is paramount. Full recovery, including return to normal activities, typically takes 4-8 weeks, though the long-term immune implications require lifelong management.

\section*{Surgical Findings \& Pathology}
During splenectomy, the surgeon documents key intraoperative findings, which may include the size and weight of the spleen, the presence of rupture or lacerations (in trauma), the extent of splenomegaly, the presence of adhesions, the appearance of the splenic capsule and parenchyma, and the identification and removal of any accessory spleens. Any abnormalities of adjacent organs, such as the pancreas or stomach, are also noted.

The excised spleen is sent for pathological examination, which is critical for confirming the underlying diagnosis. The pathology report will describe the gross appearance of the spleen (size, weight, presence of cysts, tumors, or hemorrhage) and the microscopic findings. For hematologic disorders, it will confirm the presence of specific cellular abnormalities, such as increased red blood cell destruction, lymphoid hyperplasia, or evidence of myeloproliferative disease. For oncologic indications, it will identify the type and grade of tumor (e.g., lymphoma, angiosarcoma), assess tumor margins, and determine the extent of splenic involvement. The pathologist will also confirm the presence or absence of accessory spleens in the resected tissue. This report guides further management and prognosis.

\section*{Expected Postoperative Course}
A normal, successful postoperative course following splenectomy is characterized by an uneventful recovery and the resolution of the symptoms or condition for which the surgery was performed. Patients can expect a gradual decrease in incisional pain, which is typically well-controlled with oral analgesics within a few days. Wound healing should progress without signs of infection, such as redness, swelling, or purulent discharge. Bowel function usually returns within 24-72 hours, allowing for a progressive diet. For patients with hematologic disorders, an improvement in blood counts (e.g., increased platelet count in ITP, resolution of anemia in hemolytic conditions) is expected within days to weeks. Patients should experience a return of energy and the ability to resume light activities within 1-2 weeks, with full return to normal physical activity, including strenuous exercise, typically by 4-8 weeks post-surgery. The absence of fever or other signs of infection is a key indicator of a successful early course.

\section*{Postoperative Care \& Follow-up}
Post-splenectomy follow-up is essential for monitoring recovery, managing potential complications, and ensuring long-term health, particularly regarding the lifelong risk of OPSI.

\begin{itemize}
  \item \textbf{Post-operative Clinic Visits:} Typically scheduled for 1-2 weeks and 4-6 weeks after surgery to assess wound healing, discuss pathology results, monitor for complications, and address any patient concerns.
  \item \textbf{Suture/Staple Removal:} If non-absorbable sutures or staples were used, they will be removed during an early follow-up visit, usually around 7-14 days post-op.
  \item \textbf{Complete Blood Count (CBC) Monitoring:} Regular CBCs are performed to monitor platelet counts (for thrombocytosis risk) and other blood cell lines, especially in patients with underlying hematologic disorders, to assess treatment efficacy.
  \item \textbf{Activity Resumption:} Gradual return to normal activities is advised, with specific instructions on avoiding heavy lifting or strenuous exercise for several weeks (typically 4-6 weeks) to prevent incisional hernia or wound dehiscence.
  \item \textbf{Lifelong OPSI Education:} Patients receive extensive education on the signs and symptoms of overwhelming post-splenectomy infection (fever, chills, malaise, severe headache) and are instructed to seek immediate medical attention for any febrile illness, emphasizing the need to inform healthcare providers of their splenectomy status.
  \item \textbf{Prophylactic Antibiotics:} Some patients, particularly children, those with specific risk factors, or those with ongoing immunosuppression, may be prescribed long-term prophylactic antibiotics (e.g., penicillin V). All splenectomized patients should carry a "splenectomy card" or wear a medical alert bracelet and be educated on the need for prompt empiric antibiotic treatment for any fever.
  \item \textbf{Booster Vaccinations:} Periodic booster vaccinations against encapsulated bacteria (Pneumococcal, Meningococcal, H. influenzae type B) are recommended according to national and international guidelines, typically every 5 years for pneumococcal and meningococcal vaccines.
\end{itemize}
Patients should be advised to inform all healthcare providers, including dentists, of their splenectomy status to ensure appropriate prophylactic measures are taken for invasive procedures.

\section*{Epidemiology}
The incidence of splenectomy varies significantly based on geographic region, healthcare practices, and the prevalence of specific diseases. Traumatic splenic injury remains a common indication, particularly among young males involved in high-impact accidents. However, the trend towards non-operative management for hemodynamically stable splenic trauma has significantly reduced the overall number of splenectomies performed for this indication over the past few decades. Hematologic disorders, such as immune thrombocytopenia (ITP) and hereditary spherocytosis, represent a substantial proportion of elective splenectomies, with ITP being more prevalent in younger women and children. Splenectomy for malignancy, including lymphoma or metastatic disease, is less frequent but still a significant indication. The overall mortality rate for elective splenectomy is relatively low, typically ranging from 0.5\% to 3\%, but it can be significantly higher (up to 10-20\%) in emergency settings, especially for severe trauma with associated injuries or in critically ill patients. Morbidity rates range from 10\% to 30\%, with complications such as bleeding, infection, and pancreatic injury being the most common. Overwhelming Post-Splenectomy Infection (OPSI) is a rare but devastating complication, with a lifetime risk estimated between 0.5\% and 5\%, and a mortality rate of 50-70\% if it occurs, highlighting the importance of preventative measures.

\section*{Outcomes \& Prognosis}
The outcomes and prognosis following splenectomy are largely dependent on the underlying indication for the procedure. For patients undergoing splenectomy for immune thrombocytopenia (ITP), success rates, defined as complete or partial remission (platelet count >50,000/$\mu$L without additional therapy), range from 60\% to 80\%. For hereditary spherocytosis, splenectomy typically cures the hemolytic anemia, significantly improving quality of life and reducing transfusion dependence. In cases of traumatic splenic rupture, the immediate outcome is focused on survival and prevention of re-bleeding, with long-term prognosis depending on the severity of associated injuries. For oncologic indications, the outcome is tied to the stage and type of malignancy, with splenectomy contributing to local control or staging.

Generally, the prognosis for elective splenectomy is good, with most patients experiencing resolution or significant improvement of their symptoms. Recurrence of hematologic disease can occur if accessory spleens are present and not removed, necessitating reoperation in some cases. The most significant long-term prognostic factor for all splenectomized patients is the lifelong risk of OPSI, which necessitates continuous patient education, adherence to vaccination schedules, and prompt medical attention for any signs of infection. With appropriate prophylactic measures and patient vigilance, the overall functional and survival outcomes are favorable. Landmark studies, such as those informing the American Society of Hematology guidelines for ITP, consistently show splenectomy as an effective second-line treatment, significantly improving platelet counts and reducing bleeding risk in refractory cases.

\section*{References}
\begin{enumerate}
  \item Brunicardi FC, Andersen DK, Billiar TR, et al. Schwartz's Principles of Surgery. 11th ed. McGraw-Hill Education; 2019.
  \item Neunert C, Terrell DR, Arnold DM, et al. American Society of Hematology 2019 guidelines for immune thrombocytopenia. Blood Adv. 2019 Dec 10;3(23):3829-3866.
  \item UpToDate. Splenectomy: Postoperative care and complications. Wolters Kluwer; 2024.
  \item National Institute of Diabetes and Digestive and Kidney Diseases (NIDDK). Splenectomy. U.S. Department of Health and Human Services; 2023. Available from: \url{https://www.niddk.nih.gov/health-information/digestive-diseases/splenectomy}
\end{enumerate}

\clearpage